Letting applications reuse I/O data already in kernel memory requires the kernel
to retain control over buffer lifetimes and protect its own state. Copying keeps
kernel and application buffers separate at a cost in CPU time and memory bandwidth.
Linux zero-copy sends can avoid a copy by borrowing application buffers, which
the sender must leave unchanged until the kernel signals that they can be reused.
Some existing receive interfaces require specialized network hardware to give
applications direct access to payload while keeping protocol headers and unrelated
traffic outside their buffers.

We present \sys, a Linux I/O architecture that reverses the loan. The kernel
retains ownership of I/O buffers and lends access according to the application's
authority. \emph{Named mode} gives unprivileged applications checked handles to
retained objects whose backing pages remain unmapped. \emph{Mapped mode} gives
trusted applications read/write access to a dedicated I/O region. For mapped
receive, the driver supplies buffers from that region, and the kernel keeps
control state and the header copies used for parsing outside the mapping.
A device-wide mapping requires administrative authorization because it exposes
the device's traffic. Neither mode requires hardware header splitting or
per-application flow steering.

We illustrate \sys with TCP streams, UDP datagrams, and block I/O, achieving a consistent reduction in CPU time per delivered byte compared to alternatives.
% CLAIMS S1, full-speed matched comparison. The copying comparison (S6) is unmatched.
